The feasibility assertion is exactly \(\text{lem:fixed-experiment-shell-certificate}\). Fix now \(d\ge4\), \(\eta>0\), and \(1<D\le C<e^\eta\). By \(\text{prop:zero-risk-nontrivial-shell}\), one fixed finite public experiment has a nonempty exact shell at these indices and has \(N^\star(\epsilon;d,C,D,\eta)=1\) for every \(\epsilon>0\). By \(\text{prop:informative-fast-shell-all-indices}\), another fixed finite public experiment at the same indices has the displayed local lower bound. Both propositions quantify over the full interior feasible region, so they apply without the restriction \(2D-1<e^\eta\).

Thus two public experiments having identical \((d,C,D,\eta)\) indices have different minimax sample complexities: one is identically one, while the other diverges at least as \(dD\eta/\epsilon\) as \(\epsilon\downarrow0\). Hence the fixed-experiment frontier cannot be a function of those four indices alone. Finally, any lower bound asserted uniformly over all fixed public experiments is contradicted by the zero-risk experiment whenever its right-hand side exceeds one. This applies in particular to a proposed uniform \(\Omega(dD\epsilon^{-2})\) branch. This is a partial resolution: it does not decide whether a slow branch occurs for a specified informative experiment or whether a different learner attains the sharper upper rate.